\documentclass[10pt,a4paper]{amsart}
\usepackage[margin=19mm]{geometry}
\usepackage{amsmath,amssymb,mathtools}
\usepackage[scr=rsfs,cal=euler]{mathalfa}
\usepackage{xfrac}
\usepackage[unicode,hidelinks,bookmarks=false]{hyperref}
\newcommand{\ie}{i.e.\ }
\newcommand{\cf}{cf.\ }
\newcommand{\resp}{respectively\ }
\newcommand{\Subsection}{\subsection}
\providecommand{\nobreakdash}{\nobreak\hbox{-}}
\setlength{\emergencystretch}{2em}
\multlinegap0pt
\usepackage{bm}
\usepackage{bbm}
\usepackage{dsfont}
\usepackage{xspace}
\usepackage{cancel}
\usepackage{mathtools}
\usepackage{amsthm}
\makeatletter
\@ifundefined{lemma}{\newtheorem{lemma}{Lemma}}{}
\@ifundefined{theorem}{\newtheorem{theorem}{Theorem}}{}
\makeatother
\providecommand{\vx}{\ensuremath{\vec{x}}}
\renewcommand{\vec}[1]{\boldsymbol{\mathrm{#1}}}
\title[Conductance Estimation in Digraphs: Submodular Transformatio]{Conductance Estimation in Digraphs: Submodular Transformation, Lov\'asz Extension and Dinkelbach Iteration}
\author{Sihong Shao, Chuan Yang, Xinyang Ye}
\date{Source: arXiv:2506.23131v1 (CC BY 4.0)}
\begin{document}
\maketitle

\noindent\emph{Source and licence.} Conductance Estimation in Digraphs: Submodular Transformation, Lov\'asz Extension and Dinkelbach Iteration. Version arXiv:2506.23131v1 (CC BY 4.0), distributed under the Creative Commons Attribution 4.0 International licence, as recorded in the arXiv licence metadata for this exact version and shown on the abstract page https://arxiv.org/abs/2506.23131v1. Full rights notice: licenses/arxiv-2506.23131-CC-BY-4.0.txt.\par\smallskip

\noindent\textit{Editorial scope.} Counted unit: the Theorem whose source label is \texttt{them:min frac lovasz} in the article (source0.tex lines 410--419, source label \texttt{them:min frac lovasz}), delivered together with the article's own complete original proof of it (source0.tex lines 422--453, one \texttt{proof} environment of 32 lines). No mathematics is written, repaired or re-derived: statement and proof are the article's own text line for line. Required context retained uncounted: lem:min Lovasz (lines 389--393). Every definition, display and label of those lines is retained unchanged. Omitted from this package and not redistributed: 1--388, 394--409, 420--421, 454--1360. Nothing inside a retained excerpt is omitted or altered: every retained body is delivered exactly as the article's lines are written, so no omission receipt of any kind accompanies this package. Citations: the retained excerpts carry no citation command at all, so no bibliography entry of the article is transcribed and no reference is redistributed. Reference adaptation: none was needed and none was made; every reference inside the delivered unit resolves inside it, so no unresolved reference or citation is produced. Printed slips kept literal: no printed word, symbol, hypothesis or number of the counted statement or of its proof was altered. Layout and class: the article's own class is \texttt{siamart250211} with packages including lipsum, amsfonts, graphicx, epstopdf, algorithmic, amsopn, subcaption, booktabs, amsmath, amssymb, cases; the wrapper is the lane's pinned \texttt{amsart} editorial wrapper at 10pt on A4, which restates the theorem-like environments the retained text needs and adds the article's own macro definitions that the retained text needs (\texttt{\textbackslash vec}, \texttt{\textbackslash vx}), with extra wrapper packages (bm, bbm, dsfont, xspace, cancel, mathtools). The delivered numbering is the unit's own. Assets: no retained span contains a figure environment, a graphics-inclusion command, a PDF-page inclusion, a table or a TikZ picture; no image file is copied into this package, the package declares no asset dependency and no image file is redistributed with it. Licence: version arXiv:2506.23131v1 of this article is distributed under the Creative Commons Attribution 4.0 International licence, as recorded in the arXiv licence metadata for this exact version and shown on the abstract page https://arxiv.org/abs/2506.23131v1. A full rights notice is delivered at licenses/arxiv-2506.23131-CC-BY-4.0.txt. Characters: every retained excerpt is pure ASCII, so the delivered engine, XeLaTeX with its default Latin Modern text font, reports no missing character.
\begin{lemma}
\label{lem:min Lovasz}
Let $V$ be a set of size $n$ and $P(V) = \{S : S \subseteq V\}.$ Given $g,$ $h: P(V) \rightarrow [0, +\infty),$ and define $f(S) = \min\{g(S), h(S)\},$ $\forall \, S\subset V.$ $F^l(\vx),$ $G^L(\vx),$ $H^L(\vx)$ are Lov{\'a}sz extensions of $f,$ $g,$ $h,$ respectively. Then $F^L$ is upper bounded by the minimum of $G^L$ and $H^L,$ that is,
$$F^L(\vx) \leq \min\{G^L(\vx), H^L(\vx)\}.$$
\end{lemma}

\begin{theorem}\label{them:min frac lovasz}
Let $V$ be a set of size $n$ and $P(V) = \{S : S \subseteq V\}.$ Given $f_1,$ $f_2,$ $g: P(V) \rightarrow [0, +\infty),$ $g(\emptyset)=g(V)=0,$ and $\forall \, \emptyset \neq S \subsetneq V,$ $g(S)>0.$ Define $f(S) = \min\{f_1(S), f_2(S)\},$ $\forall \, S\subset V.$
$F^L,$ $F^L_1,$ $F^L_2,$ $G^L$ are Lov{\'a}sz extensions of $f,$ $f_1,$ $f_2,$ $g,$ respectively.
%The Lov{\'a}sz extensions of $f_1, f_2,g $ are $F^L_1, F^L_2,G^L$, respectively.
%Let $f(S) = \min\{f_1(S), f_2(S)\}$, $\forall S\subset V,$ and $F^L$ is the Lov{\'a}sz extension of $f.$
Set $\Tilde{F}^L = \min\{F_1^L, F_2^L\},$ then for any set $ S\subset V $, $ F^L(\mathbf{1}_S) = \Tilde{F}^L(\mathbf{1}_S) $.
If $f_1(V)=f_2(V)=0,$ then it holds that
$$\min_{\emptyset \neq S \subsetneq V} \frac{f(S)}{g(S)} = \inf_{\vx \in \mathbb{R}^n \backslash \{ t \vec 1\}, t \in R} \frac{F^L(\vx)}{G^L(\vx)}
 = \inf_{\vx \in \mathbb{R}^n \backslash \{ t \vec 1\}, t \in R} \frac{\Tilde{F}^L(\vx)}{G^L(\vx)}.$$
\end{theorem}

\begin{proof}

For any set $S\subset V$, 
$$F^L(\mathbf{1}_S) = f(S)=\min\{f_1(S),f_2(S)\} =\min\{ F_1^L(\mathbf{1}_S),F_2^L(\mathbf{1}_S)\}=\Tilde{F}^L(\mathbf{1}_S),$$
which proves $F^L(\mathbf{1}_S) =\Tilde{F}^L(\mathbf{1}_S).$
For any $\vx \in \mathbb{R}^n \backslash \{t \vec 1\}, t \in R,$ the corresponding set $V_{\sigma(i)}$ $(1\leq i \leq n-1)$ is not empty and is a proper subset of $V,$ which leads to $g(V_{\sigma(i)})>0,$ $\forall \, 1\leq i \leq n-1$ and $G^L(\vx)>0.$
Given that $f(V)=\min\{f_1(V),f_2(V)\}= g(V)=0,$  for any $\vx \in \mathbb{R}^n \backslash \{t \vec 1\}, t \in R,$ we have
\begin{align*}
F^L(\vx)&=\sum_{i=1}^{n-1} f(V_{\sigma(i)})\left(x_{\sigma(i+1)}-x_{\sigma(i)}\right).\\
&=\sum_{i=1}^{n-1}\frac{f(V_{\sigma(i)})}{g(V_{\sigma(i)})}g(V_{\sigma(i)})\left(x_{\sigma(i+1)}-x_{\sigma(i)}\right) \\
&\geq\min_{j=1,...,n-1}\frac{f(V_{\sigma(j)})}{g(V_{\sigma(j)})}\left(\sum_{i=1}^{n-1}g(V_{\sigma(i)}) (x_{\sigma(i+1)}-x_{\sigma(i)})\right)\\
&=\min_{j=1,...,n-1}\frac{f(V_{\sigma(j)})}{g(V_{\sigma(j)})} G^L(\vx).
\end{align*}
By Lemma \ref{lem:min Lovasz},
\begin{equation*}
\frac{\Tilde{F}^L(\vx)}{G^L(\vx)}\geq 
\frac{F^L(\vx)}{G^L(\vx)}\geq \min\limits_{j=1,...,n-1} \frac{f(V_{\sigma(j)})}{g(V_{\sigma(j)})}.
\end{equation*}
On the one hand,
$$\inf_{\vx \in \mathbb{R}^n \backslash \{ t \vec 1\}, t \in R}\ \frac{\Tilde{F}^L(\vx)}{G^L(\vx)}\geq
\inf_{\vx \in \mathbb{R}^n \backslash \{ t \vec 1\}, t \in R}\frac{F^L(\vx)}{G^L(\vx)}
\geq\min\limits_{j=1,...,n-1} \frac{f(V_{\sigma(j)})}{g(V_{\sigma(j)})}
\geq\min_{\emptyset \neq S \subsetneq V}\frac{f(S)}{g(S)}. $$
On the other hand, we have $\Tilde{F}^L(\mathbf{1}_S)=F^L(\mathbf{1}_S)$ for any set $S\subset V$, thus leading to
$$\min_{\emptyset \neq S \subsetneq V}\frac{f(S)}{ g(S)}
=\min_{\emptyset \neq S \subsetneq V}\frac{F^L(\mathbf{1}_S)}{G^L(\mathbf{1}_S)}
=\min_{\emptyset \neq S \subsetneq V}\frac{\Tilde{F}^L(\mathbf{1}_S)}{G^L(\mathbf{1}_S)}
\geq\inf_{\vx \in \mathbb{R}^n \backslash \{ t \vec 1\}, t \in R}\frac{\Tilde{F}^L(\vx)}{G^L(\vx)}.$$
Therefore,
$$\min_{\emptyset \neq S \subsetneq V}\frac{f(S)}{ g(S)}
=\inf_{\vx \in \mathbb{R}^n \backslash \{ t \vec 1\}, t \in R}\frac{\Tilde{F}^L(\vx)}{G^L(\vx)}. $$
\end{proof}

\end{document}
